% ===================================================================
\section{Predation (EE05)}
% ===================================================================

\subsection{Definition and types}
\textbf{Predation}: the process by which an organism eats all or part of another organism. Reasons to study it:
economic (pests, harvests), ecological (community organisation, conservation), evolutionary (it shapes the traits of
both sides). Types:
\begin{itemize}
\item \textbf{herbivory}: animals eat plants, often partially;
\item \textbf{carnivory}: an animal \emph{or plant} eats animals (a pitcher plant counts);
\item \textbf{parasitoidism}: eggs laid on or near a host that the offspring later kill;
\item \textbf{parasitism}: living on or in a host and feeding from it without (usually) killing it;
\item \textbf{cannibalism}: predator and prey of the same species.
\end{itemize}
Unlike competition, which divides a supply, predation \emph{transfers} biomass and kills the loser, so predator and
prey numbers feed back on each other: more prey feed more predators, more predators remove more prey.

\subsection{The hare--lynx cycle: intrinsic or extrinsic?}
Hudson's Bay Company fur returns (Elton 1924), $\sim$1850--1930: snowshoe hare and lynx cycle with a period of about
10 years, the lynx following the hare. Are the cycles \textbf{intrinsic} (generated by the interacting populations:
food, predators, disease) or \textbf{extrinsic} (weather, sunspots)?

Background (Molles): the sunspot hypothesis failed (the cycles do not match); ``overpopulation'' hypotheses
(disease, stress, starvation) and ``predators drive it'' were each insufficient alone. \textbf{Krebs et al.\ (1995)},
8 years, 1-km$^2$ blocks: excluding predators doubled hare density, adding food tripled it, both together raised it
\textbf{eleven-fold}.

\begin{chain}[: what $\times2$, $\times3$, $\times11$ means]
If the two factors acted independently and additively, the joint effect would be $1+1+2=4$; independently and
multiplicatively, $2\times3=6$. The observed 11 exceeds both: the factors \textbf{interact synergistically}. A
well-fed population converts lower predation into growth, and a predator-free one converts extra food into
numbers. So no single-factor hypothesis can work: the cycle is an interaction of \textbf{three trophic levels}
(plants, hares, predators).
\end{chain}

\subsection{The Lotka--Volterra predator--prey model}
\[
\frac{\dd H}{\dd t}=rH-pHP,\qquad \frac{\dd P}{\dd t}=apHP-mP,\qquad
H^*=\frac{m}{ap},\quad P^*=\frac rp .
\]
Prey grow exponentially and are removed in proportion to encounters $HP$ (capture coefficient $p$); predators die at
rate $m$ without prey and convert captured prey into offspring with efficiency $a$.

\begin{chain}[: reading the equilibrium]
\begin{enumerate}
\item The \emph{prey} stop changing when predators are at $P^*=r/p$: just enough to eat the prey's growth.
\item The \emph{predators} stop changing when prey are at $H^*=m/(ap)$: just enough to feed them.
\item So the prey level is set by the \emph{predator's} traits ($m$, $a$, $p$), and the predator level by the
\emph{prey's} growth rate $r$. Doubling $r$ doubles the predators and leaves prey unchanged.
\end{enumerate}
\end{chain}

\begin{chain}[: why the predator lags]
Predators increase whenever $H>H^*$, so they keep increasing while prey are abundant, and peak only when the
declining prey fall back through $H^*$. Prey increase whenever $P<P^*$, so their rise is timed by the predators' crash.
Each species' turning points are set by the \emph{other} crossing its isocline: prey peak first, predators a quarter
cycle later (for small cycles).
\end{chain}

\textbf{Neutral cycles.} Linearising gives eigenvalues $\pm i\sqrt{rm}$ (small-cycle period $2\pi/\sqrt{rm}$), and the
quantity $V=apH-m\ln H+pP-r\ln P$ is conserved along every solution, so each orbit is a closed curve. The model
``only admits oscillatory behaviour''.

\textbf{Assumptions} (slide 7): (1) prey growth limited only by predation; (2) a specialist predator that needs this
prey; (3) a constant environment; (4) random encounters in a homogeneous environment.

\begin{trap}
The L--V cycles are \textbf{neutral and structurally unstable}. A shock moves the system to another orbit, where it
stays (no preferred amplitude), and any small extra term in the equations turns the closed orbits into spirals.
Real systems always have such terms, so the model's perfect cycles cannot be the explanation of real cycles.
\end{trap}

\textbf{Volterra's principle} (background): killing prey and predators at the same extra rate $\delta$ gives
$r\to r-\delta$, $m\to m+\delta$, so $H^*$ \emph{rises} and $P^*$ \emph{falls}. A broad pesticide that also kills natural
enemies can make a pest more abundant.

\subsection{Density-dependent prey, and the puzzle}
With logistic prey, $\dot H=rH(1-H/K)-pHP$: the predator isocline is still $H^*=m/(ap)$, the prey isocline slopes down,
$P=(r/p)(1-H/K)$, and $P^*=(r/p)(1-H^*/K)$, positive only if $K>m/(ap)$ (otherwise predators cannot be fed and die out).
The Jacobian now has negative trace ($-rH^*/K$) and positive determinant: a \textbf{stable equilibrium}, usually a
spiral. Prey self-limitation acts like friction on an oscillator.

\begin{keyidea}[: the three-part puzzle (slide 9)]
(1) Real data show cycles. (2) The highly unrealistic L--V model shows cycles. (3) Realistic (density-dependent)
L--V models do \emph{not}. The simple model gets the observation for the wrong reason, and adding realism loses it.
``Enter lab experiments.''
\end{keyidea}

(Background: the answer depends on which realism is added. Prey self-limitation damps cycles; a \emph{saturating}
predator (Holling type II, $g=aH/(1+ahH)$, maximum $1/h$) destabilises, because a satiated predator removes a
smaller fraction of prey as prey increase. Together they give stable limit cycles of fixed amplitude once $K$ is large:
the ``paradox of enrichment''.)

\subsection{Laboratory tests: Gause and Huffaker}
\begin{center}\small
\begin{tabularx}{\textwidth}{lX}
\toprule
Gause (\emph{Paramecium} prey, \emph{Didinium} predator) & result\\\midrule
(a) simple culture, no refuge & predator eats all prey, then starves: \textbf{both extinct}\\
(b) sediment refuge for prey & predator eats exposed prey, then starves: \textbf{predator extinct}, prey recover\\
(c) one prey + one predator added every third day & \textbf{sustained oscillations}\\\bottomrule
\end{tabularx}
\end{center}
Gause's conclusion: the system would be driven to extinction if left to itself; oscillations are possible only with
constant interference from outside.

\textbf{Huffaker}: ``Gause's system was too simple and lacked environmental heterogeneity.'' Prey mite
\emph{Eotetranychus} on oranges, predatory mite \emph{Typhlodromus}; petroleum-jelly barriers; prey disperse by
ballooning on silk, predators must crawl; only $\sim$5\,\% of each orange exposed. Simple layouts failed; the most
elaborate, \textbf{120 oranges}, gave \textbf{three oscillations over about six months}.

\begin{chain}[: why patchiness rescues the system]
Within one patch the efficient predator wins and the prey die out locally. But prey disperse faster than
predators, so they reach empty patches before predators find them. The whole system persists as local extinctions
balanced by recolonisations: a \textbf{metapopulation}. In Levins' model the occupied fraction is $p^*=1-e/c$, positive
only if colonisation $c$ exceeds extinction $e$. Hence ``if the predator is efficient, the prey can survive only in a
patchy environment''. This breaks L--V assumption (4), random encounter in a homogeneous space. Gause's immigration
in (c) does the same job as Huffaker's dispersal: rescue from outside the local system.
\end{chain}

\textbf{In nature} (slide 15): a very complicated story; all kinds of dynamics are seen, because of multiple prey and
predators, spatial heterogeneity and refugia, predator efficiency changing over time, and evolution of both sides.
Wrap-up lecture: in simple homogeneous systems prey are driven extinct; real systems are heterogeneous and both sides
evolve strategies, giving a dynamic, ever-shifting equilibrium.

\subsection{Anti-predator strategies}
\begin{center}\small
\begin{tabularx}{\textwidth}{lXl}
\toprule
strategy & logic & example\\\midrule
aposematism & conspicuous signal advertising danger or distaste & yellow-winged darter\\
crypsis & resemble the background or an object & dead-leaf mantis, tawny frogmouth\\
Batesian mimicry & harmless mimic copies a defended model's signal & drone fly / honey bee\\
Müllerian mimicry & several defended species share one signal & \emph{Heliconius}\\
catalepsis & play dead & opossum\\
intimidation display & startle or bluff & eyespot moths\\
group living & vigilance, dilution, mobbing, confusion, alarm calls/pheromones & mobbing birds, aphids\\
chemical defence & toxic, smelly, distasteful & skunk, bombardier beetle\\\bottomrule
\end{tabularx}
\end{center}

\begin{chain}[: why toxins come with warning colours]
A poisoned predator does not save the prey it has already killed. The benefit of toxicity reaches the prey
\emph{population} only if predators \emph{learn} to avoid the species after a sample. Learning needs a memorable,
recognisable signal. So distastefulness and conspicuous colour evolve together.
\end{chain}

\begin{chain}[: why a Batesian mimic must be rarer than its model]
A predator attacking a signal-bearer gains $G$ from a mimic and pays $C$ for a model. If a fraction $q$ are models, the
expected gain is $(1-q)G-qC$; it attacks only if $q<G/(G+C)$. As mimics become common, the signal loses its
information and predators start attacking: the mimic is protected only while rarer than the model. In Müllerian
mimicry every bearer is defended, so each added species only \emph{shares} the cost of educating predators.
\end{chain}

\begin{numbers}
Hare--lynx period $\approx$10\,yr. Krebs: predators out $\times2$, food $\times3$, both $\times11$. L--V: $H^*=m/(ap)$,
$P^*=r/p$, period $2\pi/\sqrt{rm}$. Huffaker: 120 oranges, 3 oscillations, $\sim$6 months, 5\,\% of each orange exposed.
\end{numbers}

\begin{statements}
\begin{enumerate}[label=\arabic*.]
\item ``In the L--V model, increasing the prey's growth rate increases the equilibrium prey density.''
\F $H^*=m/(ap)$ does not involve $r$; predators ($P^*=r/p$) absorb the extra growth.
\item ``The L--V cycles return to their original amplitude after a disturbance.''
\F They are neutral: a disturbance shifts the system to a new orbit. A cycle that returns is a \emph{limit cycle}.
\item ``Adding prey density dependence to L--V makes the cycles more realistic.''
\F It damps them to a stable point: that is part (3) of the puzzle.
\item ``In Gause's refuge experiment both species went extinct.''
\F Only the predator did; the hidden prey recovered.
\item ``Huffaker's result shows that spatial heterogeneity with differential dispersal can allow predator and prey to persist together.''
\T Prey escaped to new patches ahead of predators (metapopulation dynamics).
\item ``A toxin that kills the predator protects the individual prey that was eaten.''
\F The benefit works through predator learning, hence warning colours.
\item ``Müllerian mimics are parasites on each other's signal.''
\F All are defended; sharing the signal lowers each species' losses to predator education. Batesian mimics are the ``parasites''.
\item ``A prey trait that reduces detection by a predator shows predator--prey coevolution.''
\F Coevolution needs reciprocal change in \emph{both} parties; a one-sided adaptation is not evidence of it (2022 mid-sem Q3).
\item ``Food and predators each affect hare numbers independently.''
\F $11>2\times3$: they interact.
\end{enumerate}
\end{statements}

\begin{longans}[: ``Real predator--prey systems cycle, and so does the L--V model; yet we do not believe the model. Explain.'']
(1) Model, equilibrium, neutral closed orbits, predator lag; assumptions. (2) Structural instability: neutral orbits,
any extra term gives spirals; with logistic prey the cycles damp out. (3) The three-part puzzle. (4) Lab: Gause
(extinction unless immigration), Huffaker (patchiness and dispersal $\Rightarrow$ metapopulation persistence).
(5) Field: Krebs, three trophic levels interacting ($\times11$). (6) Nature: multiple species, heterogeneity, changing
efficiency, evolution of defences. Conclusion: cycles are real, but their cause is not the mechanism of the basic model.
\end{longans}
